\section{The Shape Diameter Function as an anatomical prior}

A separate line of investigation asked whether a local-thickness descriptor could
support anatomical assignment more robustly than surface-based rules. The
\term{shape diameter function}~\cite{sdf} assigns every surface point a
value related to the local thickness of the solid behind it: the value $\mathrm{SDF}(x)$ is obtained by casting rays into the solid
behind a surface point and taking a robust average of chord lengths, so thin
structures (legs, antennae) take small values, the bulky body takes large
ones, independently of pose; and, unlike the model's own vertex
correspondence, it is computable directly on a raw, unregistered scan. It
proved reliable for one purpose and unreliable for another. It reproduces
well across real scans (correlation $r\approx0.77$--$0.94$, $87$--$95\%$
agreement, under $1.3$\,s per specimen) and separates body from appendage on
the template with high accuracy (AUC $\approx0.973$, balanced accuracy
$\approx95.1\%$). But thresholding the field to separate \emph{individual}
appendage instances (one leg from its neighbour, left mandible from right)
fails on real scans: only $2$ of $12$ specimens had all four
expert-identified landmarks fall into distinct connected components, with
the two mandibles merged into one component in $9$ of $12$. The descriptor
is retained as a reliable, cheap semantic signal for \emph{which class of
structure} a region belongs to; it is not a mechanism for separating
individual instances of a repeated structure, because connected components cannot separate two
structures that are physically fused in the scan (laterality is not encoded in local
thickness). This is the instance-level correspondence problem this report's other sections address by different
means.
